%------------------------------------------------------------------------------%
\documentclass[fleqn,utf8,aspectratio=169,12pt]{beamer}

\usepackage{integracao_e_series}

\title{Séries Numéricas -- Teste da Comparação}

%------------------------------------------------------------------------------%
\begin{document}

\addtitle

%------------------------------------------------------------------------------%
\section{Teste da Comparação}

%------------------------------------------------------------------------------%
\definecolor{c1}{rgb}{0,.7,0}
\definecolor{c2}{rgb}{0,0,.7}
\definecolor{c3}{rgb}{.7,0,0}
\begin{frame}{Teste da Comparação}
  Sejam três séries com termos não negativos
  \[
    {\color{c1} \sum_{n=1}^{\infty} p_n} \qquad
    {\color{c2} \sum_{n=1}^{\infty} a_n} \qquad
    {\color{c3} \sum_{n=1}^{\infty} q_n}
  \]
  com \( \quad {\color{c1}p_n} \;\leq\; {\color{c2}a_n} \;\leq\; {\color{c3}q_n} \quad \) para todo $n>N$

  \pause
  \bigskip
  \begin{enumerate}
    \item se \tabto{8mm} \(\color{c3} \sum q_n \) convergir \tabto{40mm} então  \(\color{c2} \sum a_n \) converge \\[6mm]\pause
    \item se \tabto{8mm} \(\color{c1} \sum p_n \) divergir  \tabto{40mm} então  \(\color{c2} \sum a_n \) diverge
  \end{enumerate}
\end{frame}

%------------------------------------------------------------------------------%
\input{ex/G-05-Teste_comparacao-ex-1}

%------------------------------------------------------------------------------%
\section{Teste da Comparação no Limite}

%------------------------------------------------------------------------------%
\begin{frame}{Teste da Comparação no Limite}
  Sejam as séries
  \(\; \sum {\color{c1} a_n} \;\) e
  \(\; \sum {\color{c3} b_n} \;\)
  com $\;{\color{c1} a_n} > 0\;$ e $\;{\color{c3} b_n}>0\;$ para todo $n>N$
  \pause
  \vspace{\baselineskip}
  \begin{enumerate}
    \item
    \emph{Se} \(\;\lim \frac{\;{\color{c1} a_n}\;}{{\color{c3} b_n}} = c > 0 \) \;\;
    \emph{então}\; ambas convergem ou ambas divergem

    \pause
    \vfill
    \item
    \emph{Se} \(\;\sum {\color{c3} b_n}\) converge
    \tabto{36mm} \emph{e} \(\;\lim \frac{\;{\color{c1} a_n}\;}{{\color{c3} b_n}} = 0\)
    \tabto{69mm} \emph{então} \(\;\sum {\color{c1} a_n} \) também converge

    \pause
    \vfill
    \item
    \emph{Se} \(\;\sum {\color{c3} b_n}\) diverge
    \tabto{36mm} \emph{e} \(\;\lim \frac{\;{\color{c1} a_n}\;}{{\color{c3} b_n}} = \infty\)
    \tabto{69mm} \emph{então} \(\;\sum {\color{c1} a_n} \) também diverge

    \vfill
  \end{enumerate}
\end{frame}

%------------------------------------------------------------------------------%
\input{ex/G-05-Teste_comparacao-ex-2}

%------------------------------------------------------------------------------%
\section{Exemplos}

\input{ex/G-05-Teste_comparacao-ex-3}
\input{ex/G-05-Teste_comparacao-ex-4}
\input{ex/G-05-Teste_comparacao-ex-5}
\input{ex/G-05-Teste_comparacao-ex-6}
\input{ex/G-05-Teste_comparacao-ex-7}

%------------------------------------------------------------------------------%
\section{Lista Mínima}

%------------------------------------------------------------------------------%
\begin{frame}{Lista Mínima}
  Estudar a Seção 6.5 da Apostila

  \vspace{\baselineskip}
  Exercícios: 4, 7

  \vfill
  Atenção: A prova é baseada no livro, não nas apresentações
\end{frame}

%------------------------------------------------------------------------------%
\end{document}
%------------------------------------------------------------------------------%
